\documentclass[12pt]{amsart}
\usepackage{amsmath, amssymb, amsthm}
\usepackage{tikz-cd}
\usepackage{geometry}
\geometry{margin=1in}

\newcommand{\cP}{\mathcal{P}}
\newcommand{\KG}{K_G}
\newcommand{\KH}{K_H}
\newcommand{\KGH}{K_{G/H}}
\newcommand{\mdls}{\mathfrak{m}}

\newtheorem{prop}{Proposition}
\newtheorem{thm}{Theorem}
\newtheorem{lem}{Lemma}
\newtheorem{cor}{Corollary}
\theoremstyle{remark}
\newtheorem{rmk}{Remark}

\title{The Master Diagram and Unramifiedness of $\cP^{[d]_G}$ at $\eta_C$}
\author{}
\date{}

\begin{document}
\maketitle

\section{Setup and notation}

Let $\cP \to U$ be a $G$-torsor over $U = C \setminus \mdls$, where $G$ is a
finite abelian group, $H \trianglelefteq G$, and $C$ is a smooth projective
curve. Fix $p_\infty \in \mdls$ and an integer $d$ bounding the ramification
of $\cP \to U$ at $p_\infty$.

Recall the four quotient schemes
\[
\cP^{[d]_H},\qquad \cP^{[d]_G},\qquad (\cP_{G/H})^{(d)},\qquad (\cP_{G/H})^{[d]},
\]
together with the canonical Cartesian square
\[
\begin{tikzcd}[row sep=large, column sep=large]
\cP^{[d]_H} \ar[r,"H"] \ar[d]
  & (\cP_{G/H})^{(d)} \ar[d,"q"] \\
\cP^{[d]_G} \ar[r,"H"'] \ar[dr,"G"',bend right=15]
  & (\cP_{G/H})^{[d]} \ar[d,"G/H"] \\
  & U^{(d)},
\end{tikzcd}
\]
in which the arrows labelled $H$, $G/H$ and $G$ are torsors under the
indicated finite constant groups (in particular, \emph{\'etale}), and the
upper square is Cartesian (\Cref{prop:torsor_tower_diagram}).

The $G/H$-torsor $\cP_{G/H} \to U$ compactifies to a finite map of smooth
projective curves $f \colon C' \to C$. Pick $p'_\infty \in f^{-1}(p_\infty)$
with ramification index $e \le d$. Write
\[
f^{(d)} \colon C'^{(d)} \to C^{(d)}
\]
for the induced map of $d$-fold symmetric powers; it sends
$d \cdot p'_\infty \mapsto d \cdot p_\infty$. Let
\[
\pi_C \colon X_C \to C^{(d)},\qquad \pi_{C'} \colon X_{C'} \to C'^{(d)}
\]
be the blowups at $d \cdot p_\infty$ and $d \cdot p'_\infty$ respectively, with
exceptional divisors $E_C,E_{C'}$ and generic points $\eta_C,\eta_{C'}$. The
DVRs of these generic points are
\[
V_C := \mathcal{O}_{X_C,\eta_C} \subset K(C^{(d)}),\qquad
V_{C'} := \mathcal{O}_{X_{C'},\eta_{C'}} \subset K(C'^{(d)}).
\]
The local relation $t = u\cdot s^{e}$ between uniformizers at
$p_\infty,p'_\infty$ propagates to symmetric powers, giving
\[
v_{C'}(x) = e\cdot v_C(x) \quad\text{for all }x \in K(C^{(d)})^{\times},
\]
i.e.\ $V_{C'}/V_C$ is an extension of DVRs of ramification index exactly $e$.

\section{The master diagram}\label{sec:master}

The torsor tower over the open base $U^{(d)}$ embeds into the projective
geometry of $C^{(d)}$ and $C'^{(d)}$, which in turn admit canonical blowups
at the points at infinity. We assemble all of this into a single commutative
diagram.

\[
\begin{tikzcd}[row sep=2.4em, column sep=2.2em]
\cP^{[d]_H} \ar[r,"H"] \ar[d]
  & (\cP_{G/H})^{(d)} \ar[d,"q"] \ar[r,hookrightarrow,"\iota'"]
  & C'^{(d)} \ar[dd,"f^{(d)}"]
  & X_{C'} \ar[l,"\pi_{C'}"'] \ar[dd,dashed,"f_X"] \\
\cP^{[d]_G} \ar[r,"H"'] \ar[dr,"G"',bend right=10]
  & (\cP_{G/H})^{[d]} \ar[d,"G/H"]
  & & \\
  & U^{(d)} \ar[r,hookrightarrow,"\iota"']
  & C^{(d)}
  & X_C \ar[l,"\pi_C"]
\end{tikzcd}
\]

Reading the diagram:
\begin{itemize}
  \item The \textbf{left block} (columns 1--2) is the torsor tower of
        \Cref{prop:torsor_tower_diagram}: $H$-torsors horizontally, the
        $G/H$-torsor on the lower right, and the composite $G$-torsor
        $\cP^{[d]_G} \to U^{(d)}$.
  \item The \textbf{middle block} (column 3) records the compactifications.
        The map $\iota$ is the canonical open immersion
        $U^{(d)} \hookrightarrow C^{(d)}$. The map $\iota'$ is the open
        immersion $(\cP_{G/H})^{(d)} \hookrightarrow C'^{(d)}$ obtained from
        compactifying $\cP_{G/H} \subset C'$. The vertical $f^{(d)}$ is the
        finite morphism of symmetric powers; restricting along $\iota'$
        recovers the canonical map $(\cP_{G/H})^{(d)} \to U^{(d)}$, which is
        the composite of $q$ with the $G/H$-torsor.
  \item The \textbf{right block} (column 4) records the blowups at the
        diagonal points at infinity. The dashed arrow
        $f_X \colon X_{C'} \dashrightarrow X_C$ is the rational lift of
        $f^{(d)}$, defined on the strict transform; on generic points of
        the exceptional divisors it sends $\eta_{C'} \mapsto \eta_C$ and the
        induced extension of DVRs $V_C \hookrightarrow V_{C'}$ has
        ramification index $e$.
\end{itemize}

\section{Unramifiedness of $\cP^{[d]_G}$ at $\eta_C$}

\begin{thm}\label{thm:main}
Suppose
\begin{enumerate}
  \item[\textnormal{(H1)}] $(\cP_{G/H})^{[d]} \to U^{(d)}$ is unramified at
        $\eta_C$; and
  \item[\textnormal{(H2)}] $\cP^{[d]_H} \to (\cP_{G/H})^{(d)}$ is unramified
        at $\eta_{C'}$.
\end{enumerate}
Then $\cP^{[d]_G} \to U^{(d)}$ is unramified at $\eta_C$.
\end{thm}

\begin{proof}
By the tower
\[
\cP^{[d]_G} \xrightarrow{\;H\;} (\cP_{G/H})^{[d]} \xrightarrow{\;G/H\;} U^{(d)},
\]
unramifiedness of the composite at $\eta_C$ follows from unramifiedness of
each factor at the corresponding point. The lower factor is unramified at
$\eta_C$ by hypothesis (H1). It therefore suffices to prove the following.

\medskip
\noindent\textbf{Claim.} The $H$-torsor
$\cP^{[d]_G} \to (\cP_{G/H})^{[d]}$ extends to an unramified cover at any
point $\xi$ of $(\cP_{G/H})^{[d]}$ lying above $\eta_C$.
\medskip

Fix such a $\xi$. Pass to local rings:
\[
A := \mathcal{O}_{(\cP_{G/H})^{[d]},\,\xi}\,, \qquad
B := \mathcal{O}_{\cP^{[d]_G}}\!\bigl(\text{base change}\bigr),
\]
where, more precisely, $B$ is the semilocal ring obtained from
$\cP^{[d]_G} \times_{(\cP_{G/H})^{[d]}} \operatorname{Spec} A$. Since
$\cP^{[d]_G} \to (\cP_{G/H})^{[d]}$ is finite, $A \to B$ is finite, and we
must show that it is \'etale at every maximal ideal of $B$.

By (H1), the local ring $A$ is a DVR (it is regular of dimension one,
with residue field a finite separable extension of the residue field of
$V_C$, by definition of unramifiedness at $\eta_C$). The point $\eta_{C'}$
of $C'^{(d)}$ lies above $\eta_C$ via $f^{(d)}$, and via the open immersion
$\iota'$ it lies in $(\cP_{G/H})^{(d)}$. Let $\xi'$ be the unique point of
$(\cP_{G/H})^{(d)}$ above $\xi$ corresponding to $\eta_{C'}$ (so that $\xi'$
maps to $\xi$ under $q$ and to $\eta_C$ under the composite to $U^{(d)}$).
Set
\[
A' := \mathcal{O}_{(\cP_{G/H})^{(d)},\,\xi'}\,, \qquad
B' := \mathcal{O}_{\cP^{[d]_H}}\!\bigl(\text{base change to }A'\bigr).
\]
The Cartesian square in \Cref{prop:torsor_tower_diagram} localizes to a
Cartesian square of finite ring extensions
\[
\begin{tikzcd}[row sep=large, column sep=large]
B \ar[r] & B' \\
A \ar[u] \ar[r] & A' \ar[u]
\end{tikzcd}
\qquad B' \;=\; B \otimes_A A'.
\]

\textbf{Step 1: $A \to A'$ is faithfully flat.}
Both $A$ and $A'$ are DVRs (the latter by construction, since $\xi'$
specializes to the generic point of the exceptional divisor of
$X_{C'} \to C'^{(d)}$, which is regular of codimension one). The map
$A \to A'$ is the localization of the finite morphism $q$ at the chosen
points, hence it is finite. A finite, injective, local homomorphism of DVRs
is automatically flat (free of rank equal to the degree of the residue
extension times the ramification index), and is faithfully flat because the
maximal ideal of $A$ generates a non-unit ideal in $A'$.

\textbf{Step 2: $A' \to B'$ is \'etale.}
The cover $\cP^{[d]_H} \to (\cP_{G/H})^{(d)}$ is an $H$-torsor over the open
base $U^{(d)}$ pulled back along $q$, hence \'etale wherever it is defined as
such; by hypothesis (H2) it is moreover unramified at $\eta_{C'}$, hence
\'etale at $\xi'$ (a finite, unramified, flat morphism is \'etale). Therefore
$A' \to B'$ is \'etale at every maximal ideal of $B'$.

\textbf{Step 3: \'Etale descent.}
We have the Cartesian square
\[
B' = B \otimes_A A',
\]
with $A \to A'$ faithfully flat (Step 1) and $A' \to B'$ \'etale (Step 2).
By fpqc descent of \'etale morphisms (\cite[\href{https://stacks.math.columbia.edu/tag/02L4}{Tag~02L4}]{stacks}; equivalently, \cite[Exp.~IX, Prop.~4.1]{SGA1}), \'etaleness descends along faithfully flat morphisms: if
$A' \to B \otimes_A A'$ is \'etale and $A \to A'$ is faithfully flat, then
$A \to B$ is \'etale.

Therefore $A \to B$ is \'etale, i.e.\ $\cP^{[d]_G} \to (\cP_{G/H})^{[d]}$ is
\'etale at $\xi$. Since $\xi$ was an arbitrary point above $\eta_C$, we
conclude that $\cP^{[d]_G} \to (\cP_{G/H})^{[d]}$ is unramified above
$\eta_C$.

Combining with (H1), the composite $\cP^{[d]_G} \to U^{(d)}$ is unramified
at $\eta_C$, which proves the claim.
\end{proof}

\begin{rmk}
The proof is purely descent-theoretic: the only properties of the symmetric
powers and blowups used are that the relevant local rings $A$ and $A'$ are
DVRs, that $q$ is finite, and that the Cartesian square in
\Cref{prop:torsor_tower_diagram} remains Cartesian after localization.
In particular, no Abhyankar-style hypothesis on tameness or coprimality of
$e$ with the residue characteristic is needed.
\end{rmk}

\begin{rmk}[Why descent and not multiplicativity]
Naive multiplicativity of ramification indices in the diamond
\[
V_C \subset V_1 \subset W_1 \subset W_2,\qquad V_C \subset V_1 \subset V_2 \subset W_2
\]
gives only $e_1 \cdot e(W_2/W_1) = e$, where $e_1 = e(W_1/V_1)$ is the
ramification of $\cP^{[d]_G}/(\cP_{G/H})^{[d]}$ at $\xi$ that we wish to
control. Without further input this constrains $e_1$ only by $e_1 \mid e$
(in the tame range, by Abhyankar's lemma). The Cartesian-ness of the upper
square is precisely what upgrades this to the strict equality $e_1 = 1$:
it identifies $B' = B \otimes_A A'$, which is the hypothesis of fpqc descent
of \'etaleness.
\end{rmk}

\begin{thebibliography}{9}
\bibitem{SGA1} A.~Grothendieck, \emph{Rev\^{e}tements \'Etales et Groupe
Fondamental (SGA 1)}, LNM 224, Springer, 1971.
\bibitem{stacks} The Stacks Project Authors, \emph{Stacks Project},
\url{https://stacks.math.columbia.edu}.
\end{thebibliography}

\end{document}
